\thispagestyle{empty}

\begin{center}
    {\fontsize{28}{34}\selectfont\bfseries Remerciements\par}
\end{center}

\vspace{0.3cm}

{\fontsize{12}{15}\selectfont
\setstretch{1.1}
\setlength{\parindent}{0pt}
\setlength{\parskip}{0.2cm}

\textbf{\textit{Louange à Allah, Seigneur de l'univers.}}

Nous rendons grâce à \textbf{Allah, le Tout-Puissant}, par la volonté de qui toute chose aboutit, pour la santé, la force et la patience qu'Il nous a accordées afin de mener ce travail à son terme.

Nos remerciements vont d'abord à nos deux encadrants, \textit{Dr Babacar DIOP}, encadreur interne, et \textit{M. Mamadou Lamine DIEDHIOU}, encadreur externe. Tous deux se sont montrés disponibles tout au long de ce travail. Leurs conseils pertinents, leur rigueur, leur engagement, leur patience et leur expertise ont été déterminants pour mener ce mémoire à son terme.

Nous adressons également notre gratitude au \textit{Pr Abdoulaye SENE}, directeur de l'UFR, au \textit{Dr Modou LO}, chef du Département de Mathématiques et Informatique, au \textit{Pr Pascal François Mbissane FAYE}, responsable de la licence Agro TIC, au \textit{Dr Moussa GUEYE}, au \textit{Dr Madiop DIOUF}, au \textit{Dr Pape Djiby MERGANE}, à \textit{M. Bassirou DIENE} ainsi qu'à nos enseignants vacataires, pour la qualité de la formation reçue tout au long de notre parcours.

Nous renouvelons nos remerciements au \textit{Pr Pascal F. Mb FAYE}, pour avoir mis à notre disposition le capteur 8-en-1 durant tout le cycle expérimental.

Nos remerciements vont aux \textit{membres du jury} pour l'honneur qu'ils nous font en examinant ce mémoire et pour leurs remarques, qui contribueront à l'améliorer.

Nous remercions \textit{M. Idrissa TINE}, notre aîné, pour ses conseils en agroécologie.

Nous remercions nos aînés, \textit{Thiécoura NDIAYE, Mbaye Modou LO, Moustapha DIAKHATE, Bassine MBENGUE, Aïssata BALDE et Rachy S SOW}, pour leur soutien et le partage de leur expérience. 

Nos remerciements vont à nos camarades de la \textit{cinquième promotion AgroTIC}, pour leur entraide, leur solidarité et leur bienveillance.

Nous tenons également à souligner la qualité de notre collaboration et le soutien mutuel au sein de notre binôme : dans les moments de doute, notre soutien réciproque nous a rappelé que, dans la recherche, \textbf{toute erreur est une expérience}.

Enfin, nous remercions nos \textit{familles}, et en particulier nos \textit{parents}, pour leurs prières, leur patience et leurs encouragements.

Que toutes les personnes que nous n'avons pas pu citer et qui ont compté dans notre parcours trouvent ici l'expression de notre sincère gratitude.

}